\section{Evaluation methodology: three defects in our own harness}
\label{sec:methodology}

We report these because each silently corrupted a result we believed, and each
is a mistake available to any comparable pipeline.

\subsection{Selection on a maximum over noise}
\label{sec:selection}

The shipped checkpoint was chosen by the \emph{maximum} of the composite
criterion
\begin{equation}
\texttt{composite\_gain} \;=\; \texttt{hard\_gain} + \min\!\big(0,\ \texttt{clean\_gain}\big)
\label{eq:composite-gain}
\end{equation}
over 102 in-training evaluations. Across that run \texttt{clean\_gain} had mean
\SI{-1.43}{\decibel} and standard deviation \SI{1.29}{\decibel}, with only 10 of
102 evaluations positive; the shipped checkpoint's $+0.02$ ranked 9th of 102.
Its promised hard gain of \SI{+1.80}{\decibel} measured \SI{+1.43}{\decibel} on
held-out frames, between the evaluation mean ($+1.19$) and its selected maximum,
which is exactly where an inflated in-training figure should land. Selection now
uses a trailing median over five evaluations; replayed on the same series it
selects step \num{72000} instead of \num{81000}.

We then scored step \num{72000} on the same 429 frames, and the answer is not
the clean vindication we expected. On the criterion the selector optimises it is
better: composite gain \SI{-1.36}{\decibel} against the shipped checkpoint's
\SI{-1.57}{\decibel}, and clean PU21 improves by \SI{+0.55}{\decibel}. On CVVDP
it is \emph{worse on both conditions}, $-0.052$~JOD clean and $-0.071$~JOD hard,
so the smoothed rule would have shipped a checkpoint that a perceptual metric
likes less. This is the disagreement of \cref{sec:disagreement} reappearing
inside the selection rule itself: median smoothing removes the noise it was
designed to remove, and the objective it smooths is still PU21. Smoothing a
selector does not fix choosing the wrong quantity to select on. We report the
fix and its limit together, because reporting only the first would repeat the
error this section is about.

\subsection{An evaluation that read the front of the split}
\label{sec:front-of-split}

\texttt{DataLoader(val, shuffle=False)} with \texttt{max\_batches=N} evaluates
the alphabetically \emph{first} $N$ records: 32 records across 11 scenes, every
name between \texttt{abandoned\_factory} and \texttt{blau\_river}, with 403
records and 87 scenes never measured. Since error correlates with headroom
(\cref{sec:headroom}), that slice reported a clean gain of $+0.61$ where the
benchmark measures $-3.0$. Fixed with a seeded permutation applied identically
to both conditions.

\subsection{A viewer that presented every frame upside down}
\label{sec:flip}

The default framebuffer places row 0 at the bottom and the display shader
sampled its texture coordinate unchanged. No numeric test caught it, because
they all compare the float composite, which a presentation flip leaves
untouched, and the one test that read the canvas did so through
\texttt{readPixels}, which returns rows bottom-first and cancelled the flip
against itself. Master EXR output was never affected.

\subsection{Two claims we corrected against fresh measurement}
\label{sec:corrections}

\Cref{sec:tail} has been rewritten twice against new measurements. The first
draft claimed the model invents highlights; the second claimed the penalty was
tail-carried. Both were wrong, and both were plausible. We record this beside
the harness defects because the mechanism is the same: a hypothesis that
explains the headline number is not thereby established, and in both cases the
measurement that settled it was cheap.
